% Generated from locale/tr/content/lambda-calculus/introduction/currying.tex; do not hand-edit.


\olfileid[tr]{lam}{rep}{cur} 
\olsection{Currying}

Bir $\lambd$-soyutlaması $\lambd[x][M]$ tek argümanlı bir fonksiyon temsil
eder; birden çok argüman kabul eden fonksiyonlar tanımlamak istediğimizde bu
oldukça sınırlayıcıdır. Bunu yapmanın bir yolu, $\lambd$-kalkülüsünü
ikililerin, üçlülerin vb. kurulmasına izin verecek biçimde genişletmektir;
bu durumda örneğin üç konumlu bir $\lambd[x][M]$ fonksiyonu argümanının üçlü
olmasını beklerdi. Ancak bunu \emph{Currying} yoluyla yapmak daha elverişlidir.

Bir örnek düşünelim. Bir an için $\lambd$-kalkülüsünde $+$ işlemi varmış gibi
davranalım. Toplama fonksiyonu iki konumludur, yani iki argüman alır. Fakat
bir $\lambd$-soyutlaması yalnızca tek argümanlı fonksiyon verir: sözdizim
$\lambd[(x,y)][(x+y)]$ gibi ifadelere izin vermez. Bununla birlikte
$\lambd[y][(x+y)]$ ile verilen, tek argümanı $y$'ye $x$ ekleyen tek konumlu
$f_x(y)$ fonksiyonunu ele alabiliriz. Aslında bu tek bir fonksiyon değil,
her $x$ sayısı için bir tane olmak üzere farklı ``$x$ ekle'' fonksiyonlarından
oluşan bir ailedir. Şimdi başka bir tek konumlu $g$ fonksiyonunu
$\lambd[x][f_x]$ olarak tanımlayabiliriz. $x$ argümanına uygulandığında
$g(x)$, $f_x$ fonksiyonunu döndürür; dolayısıyla değerleri başka
fonksiyonlardır. Şimdi $g$'yi $x$'e, ardından sonucu $y$'ye uygularsak
$(g(x))y=f_x(y)=x+y$ elde ederiz. Böylece tek konumlu $g$, iki konumlu
toplama fonksiyonuyla aynı işi yapabilir. ``Currying'', iki konumlu
fonksiyonları değerleri tek konumlu fonksiyonlar olan tek konumlu
fonksiyonlara dönüştürme yönteminin adıdır.

İşte $\lambd$-kalkülüsünün kendi sözdiziminde bir örnek. $f(x,y)=x$
fonksiyonunu nasıl temsil ederiz? İki argüman kabul edip birincisini döndüren
bir fonksiyon tanımlamak için $\lambd[x][\lambd[y][x]]$ yazabiliriz. Bu,
kelimesi kelimesine bir $x$ argümanı kabul edip $\lambd[y][x]$ fonksiyonunu
döndüren bir fonksiyondur. $\lambd[y][x]$ başka bir $y$ argümanı kabul eder,
fakat onu atıp her zaman $x$'i döndürür. Şimdi
$\lambd[x][\lambd[y][x]]$'i iki argümana uygulayınca ne olduğuna bakalım:
\begin{align*}
  (\lambd[x][\lambd[y][x]])MN 
  \bredone & (\lambd[y][M])N \\
  \bredone & M
\end{align*}

Genel olarak $x_1$, \dots, $x_n$ parametreli ve bir $N$ terimiyle tanımlanan
bir fonksiyon yazmak için
$\lambd[x_1][\lambd[x_2][\ldots\lambd[x_n][N]]]$ kullanabiliriz. Buna $n$
argüman uygularsak şunu elde ederiz:
\begin{multline*}
  (\lambd[x_1][\lambd[x_2][\ldots\lambd[x_n][N]]]) M_1 \dots M_n \bredone\\
  \begin{aligned}
  \bredone {} & (\Subst{(\lambd[x_2][\ldots\lambd[x_n][N]])}{M_1}{x_1}) M_2
  \dots M_n\\
   \eqs {} & (\lambd[x_2][\ldots\lambd[x_n][\Subst{N}{M_1}{x_1}]]) M_2
            \dots M_n \\
  \vdots & \\
  \bredone {} &\Subst{\Subst{N}{M_1}{x_1}\ldots}{M_n}{x_n}
  \end{aligned}
\end{multline*}
Son satır, kelimesi kelimesine fonksiyon tanımının gövdesinde $x_i$ yerine
$M_i$ koymak demektir; bir fonksiyona birden çok argüman uygularken tam da
istediğimiz budur.
